\chapter{Communication Complexity — Private Coin Complexity}
%%% generated by the blueprint pipeline from TCSlib.CommunicationComplexity.NewmanTheorem.PrivateCoinComplexity
\section{Overview}

This module defines the $\varepsilon$-error private-coin randomized communication complexity
of a two-party function $f : X \to Y \to \alpha$ and establishes its basic characterization
theorems: a bound-below criterion, a bound-above criterion, an equivalent formulation via
finite-message protocols over arbitrary finite probability spaces, and monotonicity in the
error parameter $\varepsilon$.

%%%%%%%%%%%%%%%%%%%%%%%%%%%%%%%%%%%%%%%%%%%%%%%%%%%%%%%%%%%%%%%%%%%%%%%%%%%%%%%
\section{Declarations}

\begin{definition}[Private-coin communication complexity]
\lean{CommunicationComplexity.PrivateCoin.communicationComplexity}
\statementsource{roughgarden-cc-bootcamp}{
  pdf-d54d402b16e2-p002-b002,
  pdf-d54d402b16e2-p011-b003
}
\label{CommunicationComplexity.PrivateCoin.communicationComplexity}
\leanok
\uses{CommunicationComplexity.CoinTape, CommunicationComplexity.Deterministic.Protocol.complexity, CommunicationComplexity.PrivateCoin.Protocol, CommunicationComplexity.PrivateCoin.Protocol.ApproxComputes}
The $\varepsilon$-error private-coin randomized communication complexity of $f : X \to Y \to \alpha$
is the extended natural number
\[
  R_\varepsilon(f)
  \;=\;
  \inf_{\substack{n_X, n_Y \in \mathbb{N},\\ p : \mathrm{Protocol}(\{0,1\}^{n_X}, \{0,1\}^{n_Y}, X, Y, \alpha),\\ p.\mathrm{ApproxComputes}\,f\,\varepsilon}} p.\mathrm{complexity},
\]
i.e.\ the minimum worst-case number of bits exchanged over all private-coin randomized
protocols that compute $f$ with error at most $\varepsilon$ on every input.
\end{definition}

\begin{theorem}[Characterization of private-coin complexity by upper bounds]
\lean{CommunicationComplexity.PrivateCoin.communicationComplexity_le_iff}
\label{CommunicationComplexity.PrivateCoin.communicationComplexity_le_iff}
\leanok
\uses{CommunicationComplexity.CoinTape, CommunicationComplexity.Deterministic.Protocol.complexity, CommunicationComplexity.Internal.enat_iInf_le_coe_iff, CommunicationComplexity.PrivateCoin.Protocol, CommunicationComplexity.PrivateCoin.Protocol.ApproxComputes, CommunicationComplexity.PrivateCoin.communicationComplexity}
\difficulty{1}
Let $f : X \to Y \to \alpha$ be a function, let $\varepsilon \in \bbr$ be an error
bound, and let $n \in \bbn$. Then the $\varepsilon$-error private-coin randomized
communication complexity of $f$ is at most $n$ if and only if there exist coin-tape
lengths $n_X, n_Y \in \bbn$ and a private-coin protocol $p$ with Alice's randomness
space $\{0,1\}^{n_X}$ and Bob's randomness space $\{0,1\}^{n_Y}$, inputs in $X$ and $Y$,
and outputs in $\alpha$, such that $p$ computes $f$ with error at most $\varepsilon$ on
every input and the communication complexity of $p$ is at most $n$.
\end{theorem}

\begin{theorem}[Lower bound characterization of private-coin complexity]
\lean{CommunicationComplexity.PrivateCoin.le_communicationComplexity_iff}
\label{CommunicationComplexity.PrivateCoin.le_communicationComplexity_iff}
\leanok
\uses{CommunicationComplexity.CoinTape, CommunicationComplexity.Deterministic.Protocol.complexity, CommunicationComplexity.PrivateCoin.Protocol, CommunicationComplexity.PrivateCoin.Protocol.ApproxComputes, CommunicationComplexity.PrivateCoin.communicationComplexity}
\difficulty{1}
Let $f : X \to Y \to \alpha$ be a function, let $\varepsilon \in \bbr$, and let
$R_\varepsilon(f) \in \bbn \cup \{\infty\}$ denote its $\varepsilon$-error private-coin
randomized communication complexity, namely the least worst-case number of bits
exchanged by any private-coin protocol that computes $f$ with error probability at most
$\varepsilon$ on every input. Then for every $n \in \bbn$, the inequality $n \le
R_\varepsilon(f)$ holds (in the extended natural numbers) if and only if every
private-coin protocol $p$ — over coin tapes $\{0,1\}^{n_X}$ for Alice and
$\{0,1\}^{n_Y}$ for Bob, for arbitrary lengths $n_X, n_Y \in \bbn$ — that
$\varepsilon$-computes $f$ has communication complexity at least $n$.
\end{theorem}

\begin{theorem}[Finite-message characterization of private-coin complexity]
\lean{CommunicationComplexity.PrivateCoin.communicationComplexity_le_iff_finiteMessage}
\label{CommunicationComplexity.PrivateCoin.communicationComplexity_le_iff_finiteMessage}
\leanok
\uses{CommunicationComplexity.CoinTape, CommunicationComplexity.Deterministic.FiniteMessage.Protocol.complexity, CommunicationComplexity.Deterministic.FiniteMessage.Protocol.ofProtocol_complexity, CommunicationComplexity.Deterministic.FiniteMessage.Protocol.ofProtocol_run, CommunicationComplexity.Deterministic.FiniteMessage.Protocol.toProtocol, CommunicationComplexity.Deterministic.FiniteMessage.Protocol.toProtocol_complexity, CommunicationComplexity.Deterministic.FiniteMessage.Protocol.toProtocol_run, CommunicationComplexity.Deterministic.Protocol.run, CommunicationComplexity.PrivateCoin.FiniteMessage.Protocol, CommunicationComplexity.PrivateCoin.FiniteMessage.Protocol.ApproxComputes, CommunicationComplexity.PrivateCoin.FiniteMessage.Protocol.ofProtocol, CommunicationComplexity.PrivateCoin.FiniteMessage.Protocol.rrun, CommunicationComplexity.PrivateCoin.FiniteMessage.Protocol.toProtocol, CommunicationComplexity.PrivateCoin.communicationComplexity, CommunicationComplexity.PrivateCoin.communicationComplexity_le_iff}
\difficulty{4}
Let $f : X \to Y \to \alpha$ be a function, let $\varepsilon \in \bbr$ be an error
bound, and let $n \in \bbn$. Then the $\varepsilon$-error private-coin randomized
communication complexity satisfies $R_\varepsilon(f) \le n$ if and only if there exist
coin-tape lengths $n_X, n_Y \in \bbn$ and a private-coin finite-message protocol $p$
with randomness spaces $\{0,1\}^{n_X}$ and $\{0,1\}^{n_Y}$ (over input types $X$, $Y$
and output type $\alpha$) such that $p$ computes $f$ with error at most $\varepsilon$ on
every input pair and has worst-case communication cost at most $n$. Thus the
private-coin complexity, defined via binary-message protocols, is unchanged when one
instead allows protocols whose players may send elements of arbitrary finite nonempty
message alphabets, each $\beta$-valued message being charged $\lceil \log_2 \abs{\beta}
\rceil$ bits.
\end{theorem}

\begin{theorem}[Monotonicity of communication complexity in the error]
\lean{CommunicationComplexity.PrivateCoin.communicationComplexity_mono}
\label{CommunicationComplexity.PrivateCoin.communicationComplexity_mono}
\leanok
\uses{CommunicationComplexity.PrivateCoin.communicationComplexity, CommunicationComplexity.PrivateCoin.communicationComplexity_le_iff}
\difficulty{4}
Let $f : X \to Y \to \alpha$ be a function of two arguments, and let $R_\varepsilon(f)$
denote its $\varepsilon$-error private-coin randomized communication complexity — the
least worst-case number of bits exchanged by any private-coin protocol computing $f$
with probability of error at most $\varepsilon$ on every input, taken as an extended
natural number. Then for any two error bounds $\varepsilon', \varepsilon \in \bbr$ with
$\varepsilon' \le \varepsilon$, we have $R_\varepsilon(f) \le R_{\varepsilon'}(f)$; that
is, allowing a larger error can only decrease or preserve the required communication.
\end{theorem}

\begin{theorem}[Finite-probability protocols bound private-coin complexity]
\lean{CommunicationComplexity.PrivateCoin.communicationComplexity_le_of_finiteMessage}
\label{CommunicationComplexity.PrivateCoin.communicationComplexity_le_of_finiteMessage}
\leanok
\uses{CommunicationComplexity.Deterministic.FiniteMessage.Protocol.complexity, CommunicationComplexity.FiniteProbabilitySpace, CommunicationComplexity.PrivateCoin.FiniteMessage.Protocol, CommunicationComplexity.PrivateCoin.FiniteMessage.Protocol.ApproxComputes, CommunicationComplexity.PrivateCoin.FiniteMessage.Protocol.ApproxComputes_eq_ApproxSatisfies, CommunicationComplexity.PrivateCoin.FiniteMessage.Protocol.toCoinTape, CommunicationComplexity.PrivateCoin.FiniteMessage.Protocol.toCoinTape_approxSatisfies, CommunicationComplexity.PrivateCoin.FiniteMessage.Protocol.toCoinTape_complexity, CommunicationComplexity.PrivateCoin.communicationComplexity, CommunicationComplexity.PrivateCoin.communicationComplexity_le_iff_finiteMessage}
\difficulty{4}
Let $X$, $Y$, and $\alpha$ be types, let $f : X \to Y \to \alpha$ be the function to be
computed, and let $\Omega_X$ and $\Omega_Y$ be finite probability spaces serving as the
two players' private randomness. Fix real numbers $\varepsilon'$ and $\varepsilon$ with
$\varepsilon' < \varepsilon$, and let $p$ be a finite-message private-coin protocol with
randomness spaces $\Omega_X$ and $\Omega_Y$, input types $X$ and $Y$, and output type
$\alpha$, and suppose that $p$ computes $f$ with error at most $\varepsilon'$, i.e.\ for
every input pair $(x,y)$ the probability that $p$ returns a value other than $f(x,y)$ is
at most $\varepsilon'$. Then the $\varepsilon$-error private-coin randomized
communication complexity of $f$ satisfies $R_\varepsilon(f) \le \mathrm{complexity}(p)$,
where the right-hand side is the worst-case number of bits exchanged by $p$.
\end{theorem}
